This element provides a symplectic straight-pole, bending magnet with the exact
Hamiltonian in Cartesian coordinates \cite{Borland-LS356}.
The quadrupole, sextupole, and other multipole terms are defined in Cartesian coordinates.
The magnet at present is restricted to having rectangular ends.
This is quite different from \verb|CSBEND|, where the edge angles are user-defined and
where the field expansion is in curvilinear coordinates.
Strictly speaking, \verb|CSBEND| is only valid when the dipole is built with curved,
beam-following poles.

Integration of particles in \verb|CCBEND| is very similar to what's done for
\verb|KQUAD|, \verb|KSEXT|, and \verb|KOCT|. 
The only real difference is that coordinate transformations are performed at the
entrance and exit to orient the incoming central trajectory to the straight magnet axis.
In addition, the fractional strength error is adjusted to ensure that the 
outgoing central trajectory is correct.

By default, two adjustments are made at start-up and whenever the length, angle, 
gradient, or sextupole term change:
\begin{enumerate}
\item The fractional strength error is altered to ensure the correct deflecting angle. 
      This is required because the bending field varies along the trajectory.
      By default, this affects all field 
      components together, per the usual convention in {\tt elegant}. To restrict
      the strength change to the dipole term, set \verb|COMPENSATE_KN=1|.
      To turn off this optimization, set \verb|OPTIMIZE_FSE=0|.
\item The transverse position is adjusted to center the trajectory in the magnet.
      If the sagitta is $\sigma$ and \verb|ANGLE| is positive, the initial and final $x$ 
      coordinates are $x=-\sigma/2$, while the center coordinate is $x=\sigma/2$.
      To turn off this optimization, set \verb|OPTIMIZE_DX=0|.
\end{enumerate}
One can block the re-optimization of these parameters by setting 
\verb|OPTIMIZE_FSE_ONCE| and \verb|OPTIMIZE_DX_ONCE| to 1.
Note also that the optimization is performed with all error-defining parameters
(\verb|DX|, \verb|DY|, \verb|DZ|, \verb|FSE|, \verb|ETILT|, etc.) set to zero.
However, errors that are assigned to, say, the \verb|K1| value directly would
not be recognized as such.
For this reason, assigning errors to \verb|K1| is not recommended; instead, use
the \verb|FSE_QUADRUPOLE| parameter.

Having computed the ideal trajectory through a \verb|CCBEND| element, {\tt elegant}
can suppress any errors in the trajectory. Such errors may occur due to
limited accuracy in numerical integration. It is recommended to
set \verb|REFERENCE_CORRECTION=1| to ensure that the path length is corrected.
Optionally, setting \verb|REFERENCE_CORRECTION=2| would instead correct
residual transverse trajectory errors. Using \verb|REFERENCE_CORRECTION=3| corrects
both types of error.

{\bf Edge angles and edge effects} 

The user may specify edge multipoles using the \verb|EDGE_MULTIPOLE| parameter. In addition, 
the \verb|CCBEND| element supports two fringe models, selected via the \verb|FRINGEMODEL| parameter, 
which may have a value of 0 (default) or 1.
\begin{itemize}
\item[0] --- The default edge angle treatment in \verb|CCBEND| is relatively simple, consisting of a 
vertical focusing effect with momentum dependence to all orders.
Also included are edge pseudo-sextupoles (due to the body $K_1$ term) and
pseudo-octupoles (due to the body $K_2$ term).
\item[1] --- This model is based on theoretical work and code by R. Lindberg, and includes soft-fringe
  effects via a series of fringe integrals. The integrals can be computed with the companion program
  \verb|straightDipoleFringeCalc| from a generalized gradient expansion (GGE). The  GGE can be
  created using either \verb|computeCBGGE| (for cylindrical-boundary data) or \verb|computeRBGGE|
  for (rectangular-boundary data). There is an example in the {\tt elegant} examples collection.
\end{itemize}

{\bf Multipole errors}

Multipole errors are specified for the body and edge in the same fashion as for the
\verb|KQUAD| element.
The reference is the dipole field by default, but this may be changed using the
\verb|REFERENCE_ORDER| parameter.

{\bf Radiation effects}

If \verb|SYNCH_RAD| is non-zero, classical synchrotron radiation is included in tracking.
Incoherent synchrotron radiation, when requested with {\tt ISR=1},
normally uses gaussian distributions for the excitation of the electrons.
(To exclude ISR for single-particle tracking, set \verb|ISR1PART=0|.)
Setting {\tt USE\_RAD\_DIST=1} invokes a more sophisticated algorithm that
uses correct statistics for the photon energy and number distributions.
In addition, if {\tt USE\_RAD\_DIST=1} one may also set {\tt ADD\_OPENING\_ANGLE=1},
which includes the photon angular distribution when computing the effect on 
the emitting electron.  

If \verb|SYNCH_RAD| and \verb|SR_IN_ORDINARY_MATRIX| are non-zero, classical synchrotron radiation will be included in
the ordinary matrix (e.g., for \verb|twiss_output| and \verb|matrix_output|).  Symplecticity is not assured, but the results may
be interesting nonetheless. A more rigorous approach is to use \verb|moments_output|.
\verb|SR_IN_ORDINARY_MATRIX| does not affect tracking.

{\bf Adding errors}

When adding errors, care should be taken to choose the right
parameters.  The \verb|FSE|, \verb|FSE_DIPOLE|, \verb|FSE_QUADRUPOLE|, \verb|ETILT|, \verb|EPITCH|, and
\verb|YAW|, \verb|DX|, \verb|DY|, and \verb|DZ| parameters are used for
assigning errors to the strength and alignment relative to the ideal
values given by \verb|ANGLE| and \verb|TILT|.  One can also assign 
errors to \verb|ANGLE| and \verb|TILT|, but this has a different meaning:
in this case, one is assigning errors to the survey itself.  The reference
beam path changes, so there is no orbit/trajectory error.  Note that when
adding errors to \verb|FSE|, the error is assumed to come from the power
supply, which means that multipole strengths also change.

Assigning errors to \verb|K1| is also possible, but is not the best approach, since
it changes the internal reference trajectory calculation for the element.

{\bf Splitting dipoles}

The \verb|CCBEND| element does not support splitting.
{\bf Important}: Users {\em should not} attempt to split \verb|CCBEND| elements by hand, since this
will not result in the correct geometry entering and exiting the various parts.

{\bf Matrix generation}

{\tt elegant} will use tracking to determine the transport matrix for \verb|CCBEND| elements, which
is needed for computation of twiss parameters and other operations.
This can require some time, so {\tt elegant} will cache the matrices and re-use them for
identical elements.
